\chapter{Cosa avrebbero potuto lasciare}
\label{cap:repertorio_simbolico}

Il capitolo precedente ha censito ciò che è stato trovato. Questo capitolo affronta una domanda diversa: a prescindere da quello che è stato trovato, cosa avrebbero potuto produrre le popolazioni che abitavano le Alpi venete, in base a quello che sappiamo di come vivevano, di dove vivevano, e di cosa hanno prodotto popolazioni analoghe in luoghi analoghi del mondo? La risposta non può essere una lista universale: deve essere filtrata caso per caso, periodo per periodo, applicando tre criteri distinti. Il primo è l'analogia: cosa hanno lasciato popolazioni con lo stesso tipo di organizzazione economica e sociale in altri contesti ben documentati? Il secondo è la plausibilità locale: delle cose prodotte altrove, quali erano compatibili con il clima, il substrato geologico e il paesaggio delle Alpi venete? Il terzo è la conservabilità: di ciò che era plausibile produrre, cosa ha una probabilità non nulla di essere sopravvissuto fino a oggi?

Applicare questi tre filtri in sequenza non produce una lista di siti da trovare: produce una lista di aspettative ragionevoli, che è già qualcosa di più utile dell'assenza totale di aspettative con cui di fatto si affronta la ricerca in molte delle aree descritte nel capitolo precedente.

% ============================================================
\section{Il Paleolitico superiore: grotte, ripari e oggetti portatili}
% ============================================================

\subsection{Cosa è attestato in contesti analoghi}

Il Paleolitico superiore europeo (tra 45.000 e 12.000 anni prima dell'anno zero, grosso modo coincidente con le fasi fredde descritte nel quarto capitolo) è il periodo per cui il confronto mondiale è più ricco e più preciso. In Francia e Spagna, la tradizione franco-cantabrica ha prodotto centinaia di grotte con pitture parietali complesse: Chauvet (circa 36.000 anni prima dell'anno zero), Lascaux (circa 17.000), Altamira (circa 18.500). In Germania meridionale, le grotte dello Jura svevo (Hohle Fels, Vogelherd) hanno restituito le figure intagliate in avorio di mammut più antiche del mondo, tra 40.000 e 33.000 anni prima dell'anno zero. In Italia, oltre a Fumane (già discussa), la Grotta Paglicci in Puglia e la Grotta del Romito in Calabria documentano arte parietale figurativa. In Portogallo, la Valle del Côa è l'esempio più importante di arte rupestre paleolitica all'aperto in Europa: figure di animali incise su roccia scistosa lungo il fiume, datate tra 22.000 e 10.000 anni prima dell'anno zero.

In tutti questi contesti, accanto all'arte parietale nelle grotte, c'è una produzione portatile costante: figurine, bastoni traforati, elementi di ornamento in conchiglia e dente di animale, ciottoli decorati. Il Riparo Dalmeri nella Valsugana orientale, già citato nel capitolo precedente, appartiene esattamente a questa categoria: sessanta ciottoli con pitture in ocra, in contesto epigravettiano a 13.000 anni prima dell'anno zero, sono la prova che questa produzione non si fermava ai Pirenei.

\subsection{Cosa è plausibile per le Alpi venete}

Il confronto con Fumane, Riparo Tagliente (arte mobilière), Riparo Villabruna (ciottolo inciso a 14.000 anni) e Riparo Dalmeri (pietre dipinte) dimostra che tutte e tre le principali forme di produzione simbolica paleolitica erano presenti nell'area: arte parietale in grotta, arte mobilière su osso e pietra, e pittura su ciottoli. La domanda non è se queste cose venissero prodotte, ma dove e in quale quantità.

Per l'arte parietale nelle grotte, i candidati geologicamente plausibili nelle Alpi venete sono le grotte calcaree della Lessinia (dove già si scava da decenni senza trovare pitture, ma le superfici non scavate sono ancora molte) e i Colli Berici. La tradizione franco-cantabrica non si estendeva fino alle Alpi orientali con la stessa intensità del versante atlantico: le pitture di Fumane sono eccezionali proprio perché la densità di siti figurativi decresce verso est. Ma Fumane esiste, il che significa che la pratica c'era.

Le incisioni all'aperto in stile paleolitico, come quelle del Côa, sono molto meno probabili: il Côa è un caso quasi unico in Europa, e la tradizione di incidere rocce all'aperto non sembra essere stata quella dominante per i cacciatori epigravettiani delle Alpi orientali. Non è impossibile, ma non è nemmeno l'aspettativa più ragionevole per questa fase.

Quello che con più certezza veniva prodotto, e di cui quasi certamente non resterà nulla, è la decorazione su materiali organici: legno, pelle, tendine, piume. Il confronto etnografico con i cacciatori-raccoglitori storici documentati in tutto il mondo suggerisce che la maggior parte della produzione simbolica delle società di caccia sia su supporti deperibili, e che l'arte su pietra e osso rappresenti solo la punta visibile di un sistema simbolico molto più ampio.

\subsection{Ambiente e stagionalità}

I siti paleolitici delle Alpi venete mostrano un pattern ricorrente: i ripari sotto roccia e le grotte vengono occupati in inverno o nelle stagioni di transizione, mentre i siti d'alta quota (come suggerisce Dalmeri, a quota relativamente bassa ma in posizione di transito) vengono usati in estate per seguire la fauna. L'arte parietale nelle grotte è associata ai siti di occupazione stabile invernale; l'arte mobilière e i ciottoli dipinti seguono il gruppo anche in quota. Se questo schema vale per le Alpi venete come vale per il Trentino confinante, allora le grotte da cercare per l'arte parietale sono in fondovalle o a quota intermedia, non sulle creste.

% ============================================================
\section{Il Mesolitico: mobilità, quota e supporti deperibili}
% ============================================================

\subsection{Cosa è attestato in contesti analoghi}

Il Mesolitico europeo (tra 11.500 e 6.500 anni prima dell'anno zero circa, con variazioni regionali significative) è il periodo meno ricco di arte conservata in tutto l'arco alpino, e probabilmente non per caso. Le società mesolitiche erano piccoli gruppi mobili di cacciatori-raccoglitori che seguivano la fauna in cicli stagionali verticali: le basse quote in inverno, i pianori e i passi d'alta quota in estate. Questa mobilità favorisce la produzione su supporti portatili e sfavorisce la produzione su roccia fissa.

Le eccezioni documentate sono istruttive proprio perché rare. In Scandinavia, i petroglifi di Alta (Norvegia settentrionale) includono una fase mesolitica con scene di caccia. Nella penisola iberica, alcune delle incisioni della tradizione levantina sono attribuite al Mesolitico. In Italia, i ciottoli dipinti aziliani sono stati trovati in contesti del Mesolitico antico in diverse grotte del nord: sono piccoli, portatili, e hanno lasciato tracce perché si trovavano in ripari protetti. Riparo Dalmeri, con i suoi ciottoli epigravettiani, è il precedente cronologicamente più vicino a questa tradizione.

\subsection{Cosa è plausibile per le Alpi venete}

Mondeval de Sora (sepoltura con corredo a 2.150 metri) e i siti d'alta quota del Cadore e del Cansiglio documentano la presenza mesolitica intensa sulle Alpi venete. A Plan de Frea in Val Gardena, i denti infantili in un campamento d'alta quota dimostrano che non si trattava di spedizioni di soli adulti maschi ma di movimenti di gruppi familiari. Questi gruppi producevano oggetti simbolici: il corredo di Mondeval include elementi di ornamento e strumenti decorati.

Quello che è meno probabile, ma non impossibile, è che abbiano lasciato incisioni su roccia in luoghi fissi. Non perché non ne fossero capaci, ma perché la logica della mobilità non favorisce l'investimento in segni permanenti su superficie fissa. Quando lo fanno, come in Scandinavia, tendono a scegliere punti geograficamente significativi: capi lacustri, guadi fluviali, valichi. Nelle Alpi venete, i passi dolomitici frequentati dai cacciatori mesolitici (Giau, Falzarego, Valparola) sono i candidati più ragionevoli se si cercano segni di questo tipo, anche se non è ciò che ci si aspetta di trovare con più probabilità.

% ============================================================
\section{Neolitico ed Eneolitico: i segni che durano}
% ============================================================

\subsection{Cosa è attestato in contesti analoghi}

Il Neolitico europeo (approssimativamente tra 7.500 e 5.000 anni prima dell'anno zero nelle Alpi) è il periodo in cui la produzione di arte su roccia diventa sistematica e diffusa in tutto l'arco alpino. Le coppelle (cavità emisferiche poco profonde picchiettate sulla roccia) sono la forma più ubiqua: si trovano su ogni tipo di roccia, ad ogni quota, in ogni ambiente. In Valle d'Aosta ve ne sono decine di migliaia. In Trentino-Alto Adige sono documentate ovunque ci sia calcare o cristallino affiorante. In Liguria e in Piemonte costellano i ripiani rocciosi dei valichi.

Accanto alle coppelle, il Neolitico alpino ha prodotto stele antropomorfe: le statue-menhir del Trentino meridionale (Arco, Riva del Garda, Cles), le stele della Val Venosta (Lagundo, San Martino al Monte), le statue della Lunigiana, le statue sarde. Alcune di queste si trovano a pochi chilometri dal confine veneto. La Valcamonica inizia la sua produzione documentata in questo periodo. Il Monte Bego, in Liguria, ha la sua prima fase.

\subsection{Cosa è plausibile per le Alpi venete}

Le coppelle sono praticamente certe, e non averle ancora censite sistematicamente è una lacuna della ricerca, non un'assenza delle coppelle. Ogni sistema montuoso calcareo del Veneto che è stato guardato con attenzione le ha trovate: la segnalazione non sistematica sugli Euganei (trachite) e sull'Altopiano di Asiago (calcare) non è la prova che non ci siano, è la prova che nessuno le ha ancora cercate in modo comparabile a come le ha cercate qualcuno in Valle d'Aosta o in Provenza.

Le stele antropomorfe sono meno probabili come oggetti conservati (il calcare si degrada, le stele vengono riutilizzate come materiale edilizio), ma potrebbero esistere frammenti non riconosciuti in contesti di reimpiego. Le stele trentine più vicine al confine veneto (Riva del Garda, Arco) mostrano che la tradizione era presente nell'area immediatamente adiacente.

L'eneolitico aggiunge la Cultura del Vaso Campaniforme, che ha una produzione simbolica ricca (decorazioni ceramiche, pugnali in selce con foderi decorati, ornamenti in rame e oro) ma quasi nessuna arte rupestre specifica attribuita: non è una fase in cui ci si aspetta una grande espansione delle incisioni.

% ============================================================
\section{Età del bronzo: il periodo più probabile}
% ============================================================

\subsection{Cosa sappiamo già}

Il Monte Luppia dimostra direttamente che nell'età del bronzo le popolazioni delle Alpi venete producevano arte rupestre complessa: guerrieri, armi, imbarcazioni, simboli solari, coppelle. Non è una questione aperta se lo facessero: è una questione aperta dove lo facessero oltre al Garda. La Valcamonica, il Monte Bego e le incisioni di Ala documentano che la produzione rupestre dell'età del bronzo copriva tutto l'arco alpino con una densità che non ha equivalenti nei periodi precedenti.

\subsection{Cosa è plausibile altrove nel Veneto}

Ogni massiccio calcareo delle Prealpi venete con superfici rocciose esposte e frequentazione dell'età del bronzo documentata è un candidato ragionevole. La Lessinia ha frequentazione dell'età del bronzo (Sant'Anna d'Alfaedo), superfici calcaree abbondanti, e nessuna campagna di rilevamento dell'arte rupestre comparabile a quella condotta sul Monte Luppia. L'Altopiano di Asiago ha frequentazione dell'età del bronzo e superfici calcaree: stessa situazione. I versanti del Grappa verso il Brenta, con la loro morfologia di creste e ripiani, non sono stati ancora indagati sistematicamente per questo periodo.

Il dato più rilevante per orientare la ricerca è che nell'età del bronzo l'arte rupestre alpina è associata a luoghi specifici con caratteristiche ricorrenti: superfici inclinate verso il cielo o verso un bacino visivo ampio, vicine a vie di transito, in prossimità di punti d'acqua, spesso a quota intermedia tra i fondivalle abitati e le creste di pascolo estivo. Questi criteri, applicati alla morfologia delle Prealpi venete, producono una lista di candidati più ristretta e più utile di un'ispezione casuale del territorio.

% ============================================================
\section{Cosa sopravvive e cosa no}
% ============================================================

\subsection{I fattori naturali}

La conservazione di un segno su roccia dipende da tre variabili fisiche che si combinano in modo non lineare: il tipo di roccia, l'esposizione alla meteorizzazione, e il tempo trascorso.

Il calcare è il substrato più comune nelle Prealpi venete e uno dei peggiori per la conservazione delle incisioni nel lungo periodo: è solubile nell'acqua piovana acida, si desquama per effetto del gelo e disgelo ripetuto (crioclastismo), e sviluppa patine biologiche (licheni, muschi, alghe) che possono cancellare i segni più superficiali. Le incisioni a martellina (picchiettatura profonda, come al Monte Luppia e in Valcamonica) resistono meglio di quelle a graffito (incisione leggera, come alla Val d'Assa). Dopo quattromila anni, un'incisione a martellina su calcare esposto è ancora riconoscibile; un graffito sottile potrebbe essere scomparso o ridotto a una traccia illeggibile.

La dolomia è chimicamente simile al calcare ma meccanicamente più fragile ai cicli termici: i versanti dolomitici sono dominati dalle frane e dai crolli, che periodicamente rinnovano la superficie e distruggono quanto inciso. Questo non significa che non esistano incisioni dolomitiche (esistono), ma che la probabilità di conservazione è inferiore a quella su calcare compatto.

Il porfido del Lagorai è una roccia magmatica, chimicamente molto più resistente e praticamente insolubile. Se qualcuno ha inciso il porfido, quella traccia sopravvive molto meglio del calcare. Il problema è che nessuno ha ancora cercato sistematicamente incisioni sul porfido del Lagorai: è una lacuna attiva, non una risposta.

L'esposizione conta quanto il tipo di roccia. Una superficie verticale riparata da un aggetto (un'overhang) è protetta dalla pioggia diretta e dall'accumulo di acqua; si conserva incomparabilmente meglio di una superficie orizzontale esposta, su cui l'acqua stagna, il gelo agisce e la vegetazione radica. Le incisioni di Valcamonica su superfici inclinate con aggetto naturale hanno resistito cinquemila anni. Le incisioni a graffito della Val d'Assa su superfici esposte rischiano di sparire nel corso di qualche secolo.

\subsection{I fattori umani antichi}

Prima dell'età moderna, la roccia incisa veniva riutilizzata. Una superficie con coppelle dell'età del bronzo poteva ricevere nuove incisioni nell'età del ferro, poi croci medievali, poi iniziali di pastori moderni: le stratificazioni sono la norma, non l'eccezione, in tutti i complessi rupestri alpini. Questo non cancella necessariamente i segni più antichi (a Valcamonica le sovrapposizioni vengono decifrate con tecniche di rilievo tridimensionale), ma rende più difficile la lettura e talvolta fisicamente copre segni precedenti.

La costruzione di strade, terrazzamenti agricoli, edifici ha rimosso o incorporato rocce incise in tutta la pianura veneta e lungo le valli principali. Nel caso della Lessinia, la coltivazione medievale intensiva delle superfici calcaree ha probabilmente cancellato o incorporato in murature gran parte di quanto poteva esistere a quote basse.

\subsection{I fattori umani moderni}

L'estrattivismo moderno (cave di calcare, di trachite, di porfido) ha rimosso fisicamente porzioni intere di paesaggio roccioso. Le cave di trachite euganea, attive per secoli, hanno eliminato superfici che non potranno mai essere indagate. Lo stesso vale per le cave calcaree della Lessinia e per le cave di porfido della Valsugana.

La Prima Guerra Mondiale merita un paragrafo a sé, che questo libro tratta in dettaglio nel capitolo successivo e in quello sulla tafonomia: è un vettore di distruzione con caratteristiche specifiche (concentrato su determinate creste, combinato con la difficoltà di accesso successiva) che non ha equivalenti negli altri fattori considerati finora.

Il turismo e le infrastrutture turistiche moderne hanno alterato o reso inaccessibili (paradossalmente, per eccessiva frequentazione) superfici che erano rimaste intatte per millenni. Il sentiero che consolida un pianoro d'alta quota con massicciata cementizia sigilla ciò che potrebbe trovarsi sotto.

\subsection{Cosa resta, in sintesi}

La combinazione dei fattori appena descritti produce un quadro in cui le condizioni di conservazione ottimale si concentrano in un insieme di situazioni specifiche: superfici verticali calcaree con aggetto, a quota intermedia (approssimativamente tra 600 e 1.500 metri sul livello del mare), in aree non coinvolte nell'agricoltura medievale intensiva, non raggiunte dalle opere militari della Prima Guerra Mondiale, non estratte da cave moderne. Nelle Alpi venete, questo profilo esclude quasi tutta la pianura (troppo antropicizzata), le creste prealpine occidentali (guerra), e i fondivalle principali (infrastrutture). Lascia come candidati più ragionevoli: i versanti laterali della Lessinia e del Monte Baldo a quote intermedie, alcune aree dell'Altopiano di Asiago non coinvolte negli scontri del Carso interno, la fascia prealpina della provincia di Treviso (Cesen, Col Visentin), e le zone dolomitiche bellunesi già indagate con successo per la frequentazione mesolitica.

Non è una lista di promesse. È una lista di dove il rapporto tra probabilità di conservazione e probabilità che qualcosa ci fosse è più favorevole della media. La differenza tra questa lista e il silenzio delle carte che descrivono il Veneto come privo di arte rupestre è la stessa differenza che il capitolo successivo chiamerà research bias.
